\chapter[Connected neighbourhoods and topological endpoints]{Connected neighbourhoods and the topological endpoint}
\label{chap:topological-foundations}

The analytic theories considered later have different hypotheses but a
common topological task: produce small connected neighbourhoods in the
parameter set.  This chapter proves the endpoint independently of rays,
modulus, and renormalization.  It also identifies two gaps that a picture
of nested pieces can conceal: connectedness of an ambient piece need not
give connectedness of its intersection with $\M$, and a singleton
intersection of open pieces need not give a neighbourhood basis.

\section{Subspaces, connectedness, and components}

\begin{definition}[Neighbourhood in a subspace]
For $S\subseteq X$ and $x\in S$, a set $U\subseteq S$ is a neighbourhood
of $x$ if some open set $V\subseteq X$ satisfies
\[
 x\in V\cap S\subseteq U.
\]
We write $U\in\mathcal N_S(x)$.  A neighbourhood need not itself be open.
\end{definition}

\begin{definition}[Connected sets]
A space is preconnected if it cannot be written as the union of two
disjoint nonempty open subsets.  It is connected if it is preconnected
and nonempty.  A subset has either property when endowed with its
subspace topology.
\end{definition}

Lean uses \lean{IsPreconnected} for the first property and
\lean{IsConnected} for the second.  Keeping the nonemptiness condition
explicit avoids special conventions for empty pieces and absent
components.

\begin{definition}[Connected component]
For $x\in A\subseteq X$, $\operatorname{Comp}_x(A)$ is the union of all
connected subsets of $A$ containing $x$.  It is the maximal connected
subset of $A$ containing $x$.
\end{definition}

The union is connected because its members share $x$.  Thus
\[
 x\in\operatorname{Comp}_x(A)\subseteq A.
\]
If $A\subseteq B$ and $x\in A$, maximality gives
\[
 \operatorname{Comp}_x(A)\subseteq\operatorname{Comp}_x(B).
\]
Connected components are closed in their ambient subspace.  They need
not be open, even when that subspace is an open subset of a larger
space lacking local connectedness.

\begin{lemma}[Subtype and ambient connectedness]
Let $\iota:S\hookrightarrow X$ be the inclusion and let $V\subseteq S$.
Then $V$ is connected if and only if $\iota(V)$ is connected.
\end{lemma}

\begin{proof}
The inclusion is a homeomorphism from $S$ to its image with the induced
topology.  Restrict that homeomorphism to $V$; preconnectedness and
nonemptiness are preserved in both directions.
\end{proof}

This is the exact role of
\leanref{Mlc/LocalConnectivity.lean}{isConnected_subtype_val_image}.
For a parameter piece $Q\subseteq\C$, its induced subset of the
Mandelbrot subtype has image $Q\cap\M$.  Consequently, connectedness of
$Q$ alone is insufficient; the relevant condition is connectedness of
$Q\cap\M$.

\section{Local connectedness and dyadic components}

\begin{definition}[Local connectedness]
\label{def:local-connectedness}
A space $X$ is locally connected at $x$ in the neighbourhood formulation
used here if
\[
 \forall U\in\mathcal N_X(x)\quad
 \exists V\in\mathcal N_X(x),\qquad
 V\subseteq U,\quad V\text{ preconnected}.
\]
It is locally connected if this holds at every point.
\end{definition}

A neighbourhood contains its base point, so the selected $V$ is
automatically nonempty.  For the global property this formulation agrees
with the usual basis of open connected sets.
The formal definition is
\leanref{Mlc/Literature/Review.lean}{LocallyConnectedAt}; its equivalence
with \lean{LocallyConnectedSpace} is \lean{locallyConnectedAt_iff}.

\begin{definition}[Dyadic component]
For a metric space $X$, put
\[
 D_n(x)=\operatorname{Comp}_x\bigl(B_X(x,2^{-n})\bigr).
\]
\end{definition}

\begin{theorem}[The missing component condition]
\label{thm:dyadic-components}
A metric space $X$ is locally connected if and only if every $D_n(x)$
is a neighbourhood of $x$.
\end{theorem}

\begin{proof}
Assume local connectedness.  Inside the open ball choose a connected
neighbourhood of $x$.  It lies in the component $D_n(x)$, making that
component a neighbourhood.

Conversely, given $U\in\mathcal N_X(x)$, choose $n$ so large that
$B_X(x,2^{-n})\subseteq U$.  The component $D_n(x)$ is connected,
is a neighbourhood by hypothesis, and is contained in $U$.
\end{proof}

The theorem is
\leanref{Mlc/ParameterComponentApproximation.lean}{locallyConnectedSpace_iff_dyadicComponent_mem_nhds}.
The unconditional lemma \lean{dyadicComponent_connected} proves
connectedness of each component.  The difficult condition is that it be a
neighbourhood; its connectedness follows from its definition.

At an interior point of $S\subseteq\C$, sufficiently small disks lie in
$S$ and supply connected neighbourhoods directly.  Boundary points are
where this argument stops.  This explains why colouring known open
components in the Mandelbrot picture does little to address the
conjecture.

\section{Compact shrinking and neighbourhood bases}

\begin{lemma}[Compact singleton intersection]
\label{lem:compact-singleton-basis}
Let $(K_n)$ be a nested sequence of compact subsets of a Hausdorff space,
with $x\in K_n$ for every $n$.  If
\[
 \bigcap_n K_n=\{x\},
\]
then every neighbourhood $U$ of $x$ contains some $K_n$.  If in addition
each $K_n$ is a neighbourhood, the sequence is a neighbourhood basis.
\end{lemma}

\begin{proof}
Choose an open neighbourhood $V$ of $x$ contained in $U$.  If no $K_n$
were contained in $V$, all sets $K_n\setminus V$ would be nonempty,
compact, and nested.  Compactness of $K_0$ would give a point in their
intersection.  That intersection is
\[
 \left(\bigcap_n K_n\right)\setminus V=\{x\}\setminus V=\varnothing,
\]
a contradiction.  The final assertion is the definition of a basis.
\end{proof}

The general Hausdorff-space lemma is already implemented as
\leanref{Mlc/Quadratic/Complex/ParaPuzzleBasis.lean}{MLC.Quadratic.hasBasis_nhds_of_iInter_singleton}.
Its proof uses only compactness and nesting; it is valid independently
of the frozen puzzle definitions elsewhere in that file.

\begin{proposition}[Why open intersections are insufficient]
\label{prop:open-singleton-counterexample}
There are bounded open connected sets $Q_n\subseteq\C$ such that
\[
 0\in Q_{n+1}\subseteq Q_n,\qquad
 \bigcap_n Q_n=\{0\},
\]
but $(Q_n)$ is not a neighbourhood basis at $0$.
\end{proposition}

\begin{proof}
Let $a_n=1/(n+1)$ and set
\[
 Q_n=B(0,a_n)\ \cup\
      \{x+iy:0<x<2,\ 0<y<a_n\}.
\]
The disk and the rectangle are open, connected, and intersect, so their
union is connected.  The sequence is nested and contains $0$.
Every nonzero point eventually leaves the disks; a point in every
rectangle would have $0<y\le\inf_n a_n=0$, which is impossible.
Thus the intersection is $\{0\}$.  Yet $1+ia_n/2\in Q_n$ for every
$n$, so no $Q_n$ lies in $B(0,1/2)$.

The closures retain an entire interval:
\[
 \bigcap_n\overline{Q_n}=[0,2]\subset\R\subset\C.
\]
This is exactly the compactness obstruction absent from the open
intersection statement.
\end{proof}

\begin{lemma}[Proper nesting repairs the closure gap]
\label{lem:proper-nesting}
For any sequence of subsets of a topological space, if
$\overline{Q_{n+1}}\subseteq Q_n$ for every $n$, then
\[
 \bigcap_n\overline{Q_n}=\bigcap_n Q_n.
\]
\end{lemma}

\begin{proof}
One inclusion follows from $Q_n\subseteq\overline{Q_n}$.  For the other,
a point in every closure belongs to
$\overline{Q_{n+1}}\subseteq Q_n$ for each $n$.
\end{proof}

This is kernel-checked as
\leanref{Mlc/Literature/Foundations.lean}{iInter_closure_eq_of_closure_succ_subset}.
Thus bounded, properly nested parameter domains with singleton
intersection satisfy the compact closure version of the shrinking
criterion.  Without proper nesting, the closure intersection must be
proved separately.

\section{A model-independent parameter-puzzle theorem}

\begin{definition}[Connected neighbourhood basis]
A countable connected neighbourhood basis at $x\in X$ is a sequence
$(V_n)$ such that each $V_n$ is a preconnected neighbourhood of $x$ and
every neighbourhood of $x$ contains some $V_n$.
\end{definition}

Nesting is useful extra structure, but is unnecessary for this
definition or for its local-connectivity conclusion.
The Lean structure is
\leanref{Mlc/Literature/Foundations.lean}{ConnectedNeighborhoodBasis}.

\begin{theorem}[Parameter shrinking implication]
\label{thm:parameter-shrinking}
Let $c\in\M$ and let $Q_n\subseteq\C$ satisfy:
\begin{enumerate}
\item $Q_n$ is open, $c\in Q_n$, and $Q_{n+1}\subseteq Q_n$;
\item $\overline{Q_n}$ is compact for every $n$;
\item $\bigcap_n\overline{Q_n}=\{c\}$;
\item $Q_n\cap\M$ is connected for every $n$.
\end{enumerate}
Then $\M$ is locally connected at $c$.
\end{theorem}

\begin{proof}
The closures are nested compact neighbourhoods of $c$.
Lemma~\ref{lem:compact-singleton-basis} gives their cofinality in the
neighbourhood filter of the plane.  Since $Q_n\subseteq\overline{Q_n}$,
the open sets $Q_n$ also refine every ambient neighbourhood.

An open neighbourhood of $c$ in the subspace $\M$ has the form
$W\cap\M$ for some ambient open neighbourhood $W$.  Choose $n$ with
$Q_n\subseteq W$.  Then $Q_n\cap\M$ is a connected neighbourhood
of $c$ in $\M$ contained in $W\cap\M$.  This is the required local
basis.
\end{proof}

\paragraph{Lean correspondence.}
The constructor
\leanref{Mlc/Literature/Foundations.lean}{ParameterPuzzleBasis.of_compactClosure}
derives the basis field from compact closure shrinking.  The theorem
\leanref{Mlc/Literature/Foundations.lean}{localConnectivity_of_compact_parameterPieces}
states the resulting implication directly.
It imposes no Green-sublevel definition on $Q_n$.
Its hypotheses can therefore be used by a future genuine graph-cut
parapuzzle construction.

If a proof already supplies the cofinal neighbourhood-basis condition,
the shorter endpoint is
\lean{localConnectivity_of_parameterPuzzleBasis}.  The structure
\lean{ParameterPuzzleBasis} records openness, connected intersection
with $\M$, containment of the parameter, cofinality, nesting, and
singleton intersection.  Merely recording nesting and a singleton open
intersection would omit an essential hypothesis.

\section{Components of decreasing compact sets}

\begin{lemma}[Compact components]
A connected component of a compact subset of a metric space is compact.
\end{lemma}

\begin{proof}
The component is closed in the compact subspace.  Its inclusion into the
ambient space is continuous, so its ambient image is compact.
\end{proof}

\begin{theorem}[Components commute with compact decreasing limits]
\label{thm:component-intersection}
Let $(F_N)$ be a decreasing sequence of compact sets in a metric space,
and assume $x\in F_N$ for every $N$.  Then
\[
 \operatorname{Comp}_x\left(\bigcap_N F_N\right)
     =\bigcap_N\operatorname{Comp}_x(F_N).
\]
\end{theorem}

\begin{proof}
Monotonicity of components gives the left-to-right inclusion.
For the reverse inclusion, the components on the right are compact,
connected, nested, and share $x$.  Their intersection is connected:
if two disjoint relatively open sets separated it, compactness would
force that separation at some finite level, contradicting connectedness
there.  The intersection is contained in every $F_N$ and contains $x$.
Maximality of the component in $\bigcap_N F_N$ therefore gives the
reverse inclusion.
\end{proof}

The two results are
\leanref{Mlc/ParameterComponentApproximation.lean}{isCompact_component}
and
\leanref{Mlc/ParameterComponentApproximation.lean}{component_iInter_eq}.
The compactness hypothesis is doing actual work: general connected sets
can lose connectedness under intersections.

\section{Uniform outer buffers}

For a compact decreasing outer system $O_N$ with limit $S$, define
\[
 C_N(x,r)=\operatorname{Comp}_x
                  (O_N\cap\overline B(x,r)),\qquad x\in S.
\]

\begin{definition}[Uniform outer buffer]
\label{def:uniform-outer-buffer}
The pair $(S,O)$ satisfies $\mathrm{BUF}(S,O)$ if, for every $x\in S$
and $\varepsilon>0$, there are radii
\[
 0<\delta<r<\varepsilon
\]
such that
\[
 \forall N\ \exists L\ge N,\qquad
 O_L\cap\overline B(x,\delta)\subseteq C_N(x,r).
\]
\end{definition}

The radii are chosen before $N$.  The deeper stage $L$ may depend on
$N$.  This quantifier order makes a single neighbourhood survive the
entire inverse limit.

\begin{theorem}[Buffer implies local connectedness]
\label{thm:buffer-lc}
If $O_N$ are compact and decreasing, $\bigcap_N O_N=S$, and
$\mathrm{BUF}(S,O)$ holds, then $S$ is locally connected.
\end{theorem}

\begin{proof}
Fix $x\in S$ and a neighbourhood $U$ in $S$.  Choose
$\varepsilon>0$ with $S\cap B(x,\varepsilon)\subseteq U$, and obtain
$r,\delta$ from the buffer condition.  Set
\[
 D=\bigcap_N C_N(x,r).
\]
By Theorem~\ref{thm:component-intersection},
$D=\operatorname{Comp}_x(S\cap\overline B(x,r))$, so $D$ is connected
and is contained in $U$.

For $y\in S\cap B(x,\delta)$ and each $N$, choose the index $L$
provided by the buffer condition.  Since $S\subseteq O_L$, that
condition places $y$ in $C_N(x,r)$.  Hence
$S\cap B(x,\delta)\subseteq D$.  The connected set $D$ is therefore a
neighbourhood of $x$ in $S$, as required.
\end{proof}

This is
\leanref{Mlc/ParameterComponentApproximation.lean}{locallyConnectedSpace_of_uniformOuterBuffer}.
For the finite-orbit stages of $\M$, it yields the existing conditional
root.  Compactness and the exact outer limit are already proved; the
global buffer itself is an additional unproved geometric input.

\section{Small continua and the diameter modulus}

\begin{definition}[Continuum and diameter modulus]
\label{def:component-modulus}
A continuum is a nonempty compact connected metric space.
For a nonempty compact metric space $X$ and $r>0$, write
$\diam C=\sup_{u,v\in C}d(u,v)$ and define
\[
 \omega_X(r)=
 \sup_{\substack{x,y\in X\\d(x,y)<r}}
 \inf\{\diam C:C\subseteq X\text{ a continuum},\ x,y\in C\}.
\]
The infimum of the empty family is $+\infty$.
\end{definition}

\begin{theorem}[Uniform small-continuum criterion]
\label{thm:metric-lc}
For nonempty compact metric $X$, local connectedness is equivalent to:
for every $\varepsilon>0$ there is $\delta>0$ such that any
$x,y\in X$ with $d(x,y)<\delta$ lie in a continuum of diameter less
than $\varepsilon$.  Equivalently,
\[
 \lim_{r\downarrow0}\omega_X(r)=0.
\]
\end{theorem}

\begin{proof}
Assume local connectedness.  For each $a\in X$, choose a connected
neighbourhood $V_a\subseteq B_X(a,\varepsilon/4)$, and an open
neighbourhood $U_a\subseteq V_a$ of $a$.  A finite subfamily of the
$U_a$ covers $X$.  By the Lebesgue-number lemma it has a number
$\delta>0$ such that every pair at distance less than $\delta$ lies
in one member.  The closure of the corresponding $V_a$ is compact
and connected and has diameter at most $\varepsilon/2$.

Conversely, choose a ball $B_X(x,\varepsilon)$ inside a given
neighbourhood.  Apply the small-continuum condition with diameter
bound $\varepsilon$.  Every $y\in B_X(x,\delta)$ can be joined to
$x$ by a continuum contained in $B_X(x,\varepsilon)$.  Therefore
the component of that ball containing $x$ contains
$B_X(x,\delta)$ and is a connected neighbourhood.
The modulus formulation follows directly by choosing a strictly
smaller diameter bound before taking the infimum and supremum.
\end{proof}

This quantitative criterion is also used in
\lean{docs/uniform_geometric_mlc.tex}.  The proof above is included in
the mathematical account; the current Lean root uses the component
criterion rather than claiming that this particular extended-real
modulus has already been implemented.

\section{Continuous retractions}

\begin{definition}[Retraction]
Maps $i:Y\to X$ and $r:X\to Y$ form a continuous retraction if both
are continuous and $r\circ i=\id_Y$.
\end{definition}

\begin{theorem}[Local connectedness descends through a retraction]
\label{thm:retraction-lc}
If $X$ is locally connected and $i,r$ form a continuous retraction,
then $Y$ is locally connected.
\end{theorem}

\begin{proof}
Fix $y\in Y$ and $U\in\mathcal N_Y(y)$.  The preimage
$r^{-1}(U)$ is a neighbourhood of $i(y)$.  Choose a connected
neighbourhood $W$ of $i(y)$ contained in that preimage.
The image $r(W)$ is connected and lies in $U$.
Continuity of $i$ makes $i^{-1}(W)$ a neighbourhood of $y$, and
\[
 i^{-1}(W)\subseteq r(W)
\]
because $r(i(z))=z$.  Thus $r(W)$ is a connected neighbourhood
of $y$ inside $U$.
\end{proof}

The exact result is
\leanref{Mlc/RetractionLocalConnectivity.lean}{locallyConnectedSpace_of_continuous_retraction}.
It requires neither compactness nor separation assumptions.
The same file proves that convex subsets of real normed spaces are
locally connected and obtains the conditional Mandelbrot conclusion
from a supplied retraction of a convex parameter domain.
Constructing such a retraction is the remaining input; taking an
arbitrary continuous image would not supply the right-inverse
argument used in this proof.
